\documentclass[11pt]{article}
\usepackage[margin=1in]{geometry}
\usepackage[T1]{fontenc}
\usepackage[utf8]{inputenc}
\usepackage{setspace}
\usepackage{hyperref}
\hypersetup{colorlinks=true, allcolors=black}
\setlength{\parindent}{0pt}
\setlength{\parskip}{0.5em}
\setstretch{1.08}
\pagestyle{empty}

\begin{document}

\begin{center}
{\Large Teaching Statement}\\[0.3em]
Ke Cheng Hsu \\[0.2em]
{\small Duke Kunshan University faculty application (AJO \#32544)}
\end{center}

I have taught at the Hong Kong University of Science and Technology (HKUST) in several roles. From 2020 to 2025, I was a teaching assistant for undergraduate macroeconomics, the economics of uncertainty and information, and Ph.D.\ microeconomic theory, conducting weekly tutorials, holding office hours, and grading. For some courses I designed the tutorial materials myself and wrote exam questions. My average evaluation was 4.72 out of 5. I also lectured in the math camp for the MSc in Economics in 2026. I designed the material myself and covered five topics in five days: real analysis, linear algebra, statistics, optimization, and dynamic programming. Most students came from a non-math background, so I built the material step by step and kept the class encouraging. I used examples and graphs whenever they could carry the idea, such as the geometry of projection to explain OLS. Students commented that ``the teacher teaches well; let him teach more,'' and my evaluation was 4.53 out of 5. As Teaching Assistant Coordinator for the Business School from 2025 to 2026, I helped postgraduate students design teaching materials and evaluated their mock teaching.

I am prepared to teach principles and intermediate microeconomics and macroeconomics, as well as mathematics for economists. At the graduate level, I can teach Ph.D.\ microeconomic theory, information economics, applied microeconometrics, and empirical industrial organization. I can also offer an introductory course in computational methods for economics. I work in R, Python, and C++, manage projects with Git and GitHub, and use AI tools carefully in coding and debugging. I treat Git as an essential skill and teach it through a real project, since students need version control to collaborate and keep their work reproducible.

My classroom approach is student-centered and example-first. I want students to use models not just to solve equations, but to distinguish an observed pattern from the mechanism that produces it. A new restaurant on campus illustrates this. Students may treat a long line as evidence that the food is good, but the line is only a signal. As more people observe and join, the line can keep growing even when the food is bad. This example introduces herding under imperfect information. Another example asks why a government might issue consumption vouchers instead of cash. Because vouchers must be spent, more of the transfer enters current demand, and the multiplier is larger. I introduce the formula only after students understand the intuition of the underlying mechanism. I encourage students to stop me whenever something is unclear, because one student's question often helps several others who share the same doubt. If an explanation survives their questions in plain language, it usually survives future classes too.

At Duke Kunshan University I would emphasize close interaction, inclusive discussion, and applications that travel across cultures in an international classroom. In a project-based course such as Signature Work, I would give each student tailored, individual feedback at every stage, from framing a feasible question, through model and data work, to the final presentation. Teaching is most fulfilling when a student who was stuck begins to see why an outcome occurs. Finding that clearer explanation is what I enjoy most.

\end{document}
